\chapter{Généralités sur les fonctions numériques}
\textit{La notion de fonction est au cœur de toute l'analyse. Avant d'apprendre à
dériver ou à étudier des courbes, il faut savoir précisément ce qu'est une fonction,
où elle est définie, et reconnaître ses symétries et son sens de variation. Ce
chapitre pose ces fondations — elles te serviront dans tout le reste du programme.}

\section{Fonction numérique et ensemble de définition}

\begin{definition}[Fonction numérique d'une variable réelle]
Une \textbf{fonction numérique} $f$ associe à chaque réel $x$ d'une partie $D$ de
$\R$ \emph{au plus un} réel, noté $f(x)$ et appelé \textbf{image} de $x$ par $f$. On note
\[ f : D \to \R,\qquad x \mapsto f(x). \]
Si $y = f(x)$, on dit que $x$ est un \textbf{antécédent} de $y$ par $f$.
\end{definition}

\begin{remarque}
Un réel possède \emph{au plus une} image (c'est ce qui fait de $f$ une fonction),
mais un réel $y$ peut avoir \emph{plusieurs} antécédents, ou aucun. Par exemple, pour
$f(x)=x^2$, le réel $4$ a deux antécédents ($-2$ et $2$) et le réel $-1$ n'en a aucun.
\end{remarque}

\begin{definition}[Ensemble de définition]
L'\textbf{ensemble de définition} de $f$, noté $D_f$, est l'ensemble des réels $x$
pour lesquels $f(x)$ \emph{existe} (c'est-à-dire peut être calculé).
\end{definition}

\begin{aretenir}
Pour déterminer $D_f$, on traque les opérations « interdites » :
\begin{itemize}
  \item un \textbf{dénominateur} doit être non nul ;
  \item l'expression sous une \textbf{racine carrée} doit être positive ou nulle ;
  \item (plus tard : l'argument d'un $\ln$ doit être strictement positif).
\end{itemize}
$D_f$ est alors l'intersection de toutes les conditions trouvées.
\end{aretenir}

\begin{exemple}
Déterminons $D_f$ pour $f(x)=\dfrac{\sqrt{x+2}}{x-3}$.
\begin{itemize}
  \item Racine : $x+2\ge 0$, soit $x\ge -2$.
  \item Dénominateur : $x-3\neq 0$, soit $x\neq 3$.
\end{itemize}
D'où $D_f = [-2\,;\,3[\,\cup\,]3\,;\,+\infty[$.
\end{exemple}

\section{Représentation graphique}

\begin{definition}[Courbe représentative]
Le plan étant muni d'un repère $(O,\vec{\imath},\vec{\jmath})$, la \textbf{courbe
représentative} de $f$ est l'ensemble
\[ \mathcal{C}_f = \big\{\, M(x\,;\,y)\ \mid\ x\in D_f \ \text{et}\ y=f(x) \,\big\}. \]
Un point $M(x\,;\,y)$ appartient à $\mathcal{C}_f$ \textbf{si et seulement si}
$x\in D_f$ et $y=f(x)$.
\end{definition}

\begin{illus}
\begin{tikzpicture}[scale=0.95]
  \draw[->,onbline] (-3.2,0)--(3.3,0) node[right,onbmuted]{$x$};
  \draw[->,onbline] (0,-0.6)--(0,3.2) node[above,onbmuted]{$y$};
  \draw[thick,onbo,smooth,samples=80,domain=-1.7:1.7] plot (\x,{\x*\x});
  \filldraw[onbo] (1.3,1.69) circle (1.6pt) node[right,onbink]{$M(x\,;\,f(x))$};
  \draw[dashed,onbmuted] (1.3,0) node[below,onbink]{$x$} -- (1.3,1.69) -- (0,1.69) node[left,onbink]{$f(x)$};
\end{tikzpicture}
\fig{Figure — lecture graphique : à $x$ on associe l'ordonnée $f(x)$ du point de la courbe.}
\end{illus}

\section{Parité et éléments de symétrie}

\begin{definition}[Fonction paire, fonction impaire]
Soit $f$ définie sur un ensemble $D_f$ \textbf{symétrique par rapport à $0$}
(c'est-à-dire : $x\in D_f \Rightarrow -x\in D_f$).
\begin{itemize}
  \item $f$ est \textbf{paire} si, pour tout $x\in D_f$, $\ f(-x)=f(x)$ ;
  \item $f$ est \textbf{impaire} si, pour tout $x\in D_f$, $\ f(-x)=-f(x)$.
\end{itemize}
\end{definition}

\begin{propriete}[Interprétation graphique]
Dans un repère orthogonal :
\begin{itemize}
  \item si $f$ est \textbf{paire}, $\mathcal{C}_f$ admet l'\textbf{axe des ordonnées}
        $(Oy)$ pour axe de symétrie ;
  \item si $f$ est \textbf{impaire}, $\mathcal{C}_f$ admet l'\textbf{origine} $O$
        pour centre de symétrie.
\end{itemize}
\end{propriete}

\begin{illus}
\begin{tikzpicture}[scale=0.8]
  \begin{scope}
    \draw[->,onbline] (-2.4,0)--(2.4,0); \draw[->,onbline] (0,-0.4)--(0,2.6);
    \draw[thick,onbo,smooth,samples=60,domain=-1.5:1.5] plot (\x,{\x*\x});
    \node[onbmuted,below] at (0,-0.4) {paire : $f(x)=x^2$};
  \end{scope}
  \begin{scope}[xshift=6cm]
    \draw[->,onbline] (-2.4,0)--(2.4,0); \draw[->,onbline] (0,-2.2)--(0,2.2);
    \draw[thick,onbgreen,smooth,samples=60,domain=-1.35:1.35] plot (\x,{\x*\x*\x});
    \node[onbmuted,below] at (0,-2.2) {impaire : $f(x)=x^3$};
  \end{scope}
\end{tikzpicture}
\fig{Figure — symétrie axiale (fonction paire) et symétrie centrale (fonction impaire).}
\end{illus}

\begin{exemple}
Étudions la parité de $f(x)=\dfrac{x}{x^2+1}$. Son domaine $D_f=\R$ est symétrique
par rapport à $0$. Pour tout $x\in\R$ :
\[ f(-x)=\frac{-x}{(-x)^2+1}=\frac{-x}{x^2+1}=-f(x). \]
Donc $f$ est \textbf{impaire} : sa courbe est symétrique par rapport à $O$.
\end{exemple}

\begin{theoreme}[Axe et centre de symétrie (cas général)]
Soit $a,b\in\R$.
\begin{itemize}
  \item La droite d'équation $x=a$ est un \textbf{axe de symétrie} de $\mathcal{C}_f$
        si et seulement si, pour tout $x$ tel que $a+x\in D_f$ : $\ f(a-x)=f(a+x)$.
  \item Le point $\Omega(a\,;\,b)$ est un \textbf{centre de symétrie} de
        $\mathcal{C}_f$ si et seulement si, pour tout $x$ tel que $a+x\in D_f$ :
        $\ f(a-x)+f(a+x)=2b$.
\end{itemize}
\end{theoreme}

\begin{remarque}
La parité n'est que le cas particulier $a=0$ (axe $x=0$) et $\Omega=O$ (centre
$(0\,;\,0)$). Le théorème permet de détecter une symétrie ailleurs qu'à l'origine.
\end{remarque}

\section{Sens de variation}

\begin{definition}[Fonction croissante, décroissante]
Soit $f$ définie sur un intervalle $I$. On dit que $f$ est :
\begin{itemize}
  \item \textbf{croissante} sur $I$ si elle conserve l'ordre :
        $x_1<x_2 \Rightarrow f(x_1)\le f(x_2)$ (strictement croissante si $f(x_1)<f(x_2)$) ;
  \item \textbf{décroissante} sur $I$ si elle inverse l'ordre :
        $x_1<x_2 \Rightarrow f(x_1)\ge f(x_2)$ (strictement décroissante si $f(x_1)>f(x_2)$) ;
  \item \textbf{monotone} sur $I$ si elle est croissante sur $I$ ou décroissante sur $I$.
\end{itemize}
\end{definition}

\begin{aretenir}[Étudier le sens de variation par le taux de variation]
Pour deux réels distincts $x_1,x_2$ de $I$, on forme le \textbf{taux de variation}
\[ T=\frac{f(x_2)-f(x_1)}{x_2-x_1}. \]
\begin{itemize}
  \item Si $T>0$ pour tous $x_1\neq x_2$, alors $f$ est strictement croissante sur $I$ ;
  \item si $T<0$ pour tous $x_1\neq x_2$, alors $f$ est strictement décroissante sur $I$.
\end{itemize}
(En Terminale, le signe de la dérivée remplacera ce calcul.)
\end{aretenir}

\begin{exemple}
Montrons que $f(x)=x^2$ est strictement décroissante sur $]-\infty\,;\,0]$. Pour
$x_1<x_2\le 0$ :
\[ T=\frac{x_2^2-x_1^2}{x_2-x_1}=\frac{(x_2-x_1)(x_2+x_1)}{x_2-x_1}=x_1+x_2. \]
Comme $x_1<0$ et $x_2\le 0$, on a $x_1+x_2<0$, donc $T<0$ : $f$ est strictement
décroissante sur $]-\infty\,;\,0]$. (On montrerait de même qu'elle est strictement
croissante sur $[0\,;\,+\infty[$.)
\end{exemple}

\begin{definition}[Majorant, minorant, extremum]
Soit $f$ définie sur $D$ et $M,m\in\R$.
\begin{itemize}
  \item $f$ est \textbf{majorée} par $M$ si $f(x)\le M$ pour tout $x\in D$ ;
        \textbf{minorée} par $m$ si $f(x)\ge m$ pour tout $x\in D$ ; \textbf{bornée}
        si elle est à la fois majorée et minorée.
  \item $f$ admet un \textbf{maximum} $M$ en $x_0$ si $f(x_0)=M$ et $f(x)\le M$ pour
        tout $x$ ; \textbf{minimum} de façon analogue.
\end{itemize}
\end{definition}

\section{Fonctions associées et transformations de courbes}

À partir d'une fonction $f$ de courbe $\mathcal{C}_f$, on construit de nouvelles
fonctions dont la courbe se déduit de $\mathcal{C}_f$ par une transformation simple.
On note $\vec{u}$ le vecteur de translation.

\begin{propriete}[Transformations usuelles]
Soit $k>0$ et $a$ un réel.
\begin{center}\renewcommand{\arraystretch}{1.35}
\begin{tabular}{@{}l l@{}}
\toprule
\textbf{Nouvelle fonction} & \textbf{Sa courbe se déduit de $\mathcal{C}_f$ par\dots} \\
\midrule
$x\mapsto f(x)+k$   & translation de vecteur $k\,\vec{\jmath}$ (vers le haut) \\
$x\mapsto f(x)-k$   & translation de vecteur $-k\,\vec{\jmath}$ (vers le bas) \\
$x\mapsto f(x-a)$   & translation de vecteur $a\,\vec{\imath}$ (vers la droite) \\
$x\mapsto -f(x)$    & symétrie par rapport à l'axe $(Ox)$ \\
$x\mapsto f(-x)$    & symétrie par rapport à l'axe $(Oy)$ \\
$x\mapsto |f(x)|$   & on conserve la partie au-dessus de $(Ox)$, on rabat l'autre \\
\bottomrule
\end{tabular}
\end{center}
\end{propriete}

\begin{illus}
\begin{tikzpicture}[scale=0.85]
  \draw[->,onbline] (-2.6,0)--(2.8,0) node[right,onbmuted]{$x$};
  \draw[->,onbline] (0,-0.5)--(0,3.4) node[above,onbmuted]{$y$};
  \draw[thick,onbmuted,smooth,samples=60,domain=-1.5:1.5] plot (\x,{\x*\x}) node[right]{$\mathcal{C}_f$};
  \draw[thick,onbo,smooth,samples=60,domain=-1.5:1.5] plot (\x,{\x*\x+1}) node[right]{$f(x)+1$};
\end{tikzpicture}
\fig{Figure — la courbe de $x\mapsto f(x)+1$ est $\mathcal{C}_f$ translatée d'un cran vers le haut.}
\end{illus}

\begin{remarque}
La fonction $x\mapsto |f(x)|$ est très utile : elle « replie » au-dessus de l'axe des
abscisses toute la portion de courbe située en dessous. On la rencontre dès qu'un
énoncé fait intervenir une valeur absolue.
\end{remarque}

% ════════════════════════════════════════════════════════════
\section{L'essentiel à retenir}

\begin{synthese}
\textbf{Ensemble de définition.} On exclut : les valeurs annulant un dénominateur,
celles rendant négatif un radicande. $D_f$ = intersection des conditions.

\smallskip
\textbf{Appartenance à la courbe.} $M(x\,;\,y)\in\mathcal{C}_f \iff x\in D_f$ et $y=f(x)$.

\smallskip
\textbf{Parité} (sur un domaine symétrique par rapport à $0$) :
\[ f \text{ paire} \iff f(-x)=f(x)\ (\text{axe } Oy);\qquad
   f \text{ impaire} \iff f(-x)=-f(x)\ (\text{centre } O). \]
\textbf{Symétries générales} : axe $x=a$ si $f(a-x)=f(a+x)$ ; centre $\Omega(a\,;\,b)$
si $f(a-x)+f(a+x)=2b$.

\smallskip
\textbf{Sens de variation} via le taux $T=\dfrac{f(x_2)-f(x_1)}{x_2-x_1}$ :
$T>0$ (croissante), $T<0$ (décroissante).

\smallskip
\textbf{Transformations} : $f(x)+k$ (haut), $f(x-a)$ (droite), $-f(x)$ (symétrie $Ox$),
$f(-x)$ (symétrie $Oy$), $|f(x)|$ (rabattement).
\end{synthese}

% ════════════════════════════════════════════════════════════
\section{Méthodes \& savoir-faire}

\methode{Déterminer un ensemble de définition}
Repérer chaque dénominateur (le poser $\neq 0$) et chaque radicande (le poser
$\ge 0$), résoudre chaque condition, puis \emph{intersecter}.
\begin{exemple}
$f(x)=\dfrac{1}{\sqrt{4-x^2}}$. Ici le radicande est sous une racine \emph{et} au
dénominateur : il faut $4-x^2>0$ (strictement, à cause du dénominateur), soit
$-2<x<2$. Donc $D_f=\,]-2\,;\,2[$.
\end{exemple}

\methode{Étudier la parité d'une fonction}
\textbf{1.} Vérifier que $D_f$ est symétrique par rapport à $0$ (sinon, $f$ n'est ni
paire ni impaire). \textbf{2.} Calculer $f(-x)$ et le comparer à $f(x)$ et $-f(x)$.
\begin{exemple}
$f(x)=x^4-3x^2+1$ sur $\R$. On a $f(-x)=(-x)^4-3(-x)^2+1=x^4-3x^2+1=f(x)$ :
$f$ est \textbf{paire}.
\end{exemple}

\methode{Étudier le sens de variation par le taux de variation}
Former $T=\dfrac{f(x_2)-f(x_1)}{x_2-x_1}$, le simplifier (souvent en factorisant
$x_2-x_1$), puis déterminer son signe sur l'intervalle étudié.
\begin{exemple}
$f(x)=\dfrac{1}{x}$ sur $]0\,;\,+\infty[$. Pour $0<x_1<x_2$ :
\[ T=\frac{\frac1{x_2}-\frac1{x_1}}{x_2-x_1}=\frac{x_1-x_2}{x_1x_2(x_2-x_1)}=\frac{-1}{x_1x_2}<0. \]
Donc $f$ est strictement décroissante sur $]0\,;\,+\infty[$.
\end{exemple}

\methode{Démontrer qu'une droite ou un point est élément de symétrie}
Pour un axe $x=a$, prouver $f(a-x)=f(a+x)$. Pour un centre $\Omega(a\,;\,b)$, prouver
$f(a-x)+f(a+x)=2b$.
\begin{exemple}
Montrons que $\Omega(0\,;\,0)$… (cas de l'imparité, déjà vu). Plus subtil :
pour $f(x)=\dfrac{2x-1}{x-1}$, on vérifie que $\Omega(1\,;\,2)$ est centre de symétrie
en calculant $f(1-x)+f(1+x)=4=2\times 2$.
\end{exemple}

\methode{Tracer la courbe d'une fonction associée}
Partir de $\mathcal{C}_f$ connue et appliquer la transformation correspondante
(translation, symétrie, rabattement) sans recalculer point par point.
\begin{exemple}
Connaissant la courbe de $x\mapsto x^2$, celle de $x\mapsto (x-2)^2+1$ s'obtient en
translatant de $2$ vers la droite puis de $1$ vers le haut (sommet en $(2\,;\,1)$).
\end{exemple}
